\section{How This Book Should Be Read: The Change in Game}

This book begins in a place that can feel almost too simple: sentences, variables, words like ``every'' and ``some.'' That is deliberate. The early foundations of proof are not advanced because they are complicated. They are advanced because they are easy to underestimate.

Many students arrive from calculation-heavy mathematics with a strong set of useful reflexes. They can solve systems, differentiate, integrate, diagonalize a matrix, separate a differential equation, or use a numerical method to approximate an answer. Those skills matter. They are not being thrown away. But proof-based mathematics asks for a different kind of control.

The old game often asks:
\[
\text{What is the answer?}
\]
The new game asks:
\[
\text{What exactly is being claimed, and why must it be true?}
\]
That change is small in wording and large in consequences.

\begin{focusbox}
Do not measure progress through this part by pages completed. Measure it by whether you can close the file and explain the idea without leaning on the notation as a crutch.
\end{focusbox}

The purpose of this opening part is to make the new game visible. We will treat logic not as a cold symbolic machine, but as the grammar of mathematical responsibility. If a theorem says ``for every,'' you must know who chooses the object. If it says ``there exists,'' you must know what counts as producing one. If a variable appears, you must know whether it is free, fixed, arbitrary, bound, or chosen as a witness. If a sentence has no truth value yet, you must not pretend it is a theorem.

This is the level where beginners often feel silly asking questions. Do not feel silly: every proof depends on these details.

\subsection{The Reader's Contract}

Here is the contract for reading these foundations.

\begin{enumerate}
\item Read a paragraph once for the idea.
\item Read it again for the exact mathematical claim.
\item Produce your own example.
\item Produce a non-example.
\item Ask what would break if a word were changed.
\item Try to say the same thing in ordinary language, symbols, and proof language.
\end{enumerate}

The last item is important. A concept is not yours until you can move between languages. For example, consider:
\[
\forall n\in\Z,\quad n^2\geq 0.
\]
Ordinary language:

\begin{quote}
Every integer has a nonnegative square.
\end{quote}

Proof-language opening:

\begin{quote}
Let $n\in\Z$ be arbitrary. We must show $n^2\geq 0$.
\end{quote}

Symbolic form:
\[
\forall n\in\Z,\ P(n),
\qquad\text{where }P(n)\text{ is }n^2\geq0.
\]

Those are not three unrelated things. They are the same claim viewed through three lenses. Much of proof maturity is the ability to change lenses without changing the claim.

\begin{intuitionbox}
Mathematics is not made of symbols. It is made of ideas disciplined by symbols. The discipline matters. So does the idea.
\end{intuitionbox}

\subsection{The Visual Language of This Book}

As the book becomes more visual, the diagrams use a consistent language. You do
not need to memorize the colors, but the repetition should help your eye.

\begin{center}
\begin{tikzpicture}[
  every node/.style={font=\small},
  sample/.style={draw, rounded corners, minimum width=2.5cm, minimum height=0.75cm, align=center}
]
\node[sample, fill=sourceblue!10, draw=sourceblue] (s) {source / domain};
\node[sample, fill=targetgreen!10, draw=targetgreen, right=0.8cm of s] (t) {target / codomain};
\node[sample, fill=quotientorange!12, draw=quotientorange, right=0.8cm of t] (q) {quotient / gluing};
\node[sample, fill=trapred!8, draw=trapred, below=0.7cm of s] (r) {trap / failure};
\node[sample, fill=structurepurple!10, draw=structurepurple, right=0.8cm of r] (p) {structure};
\draw[-{Latex}, thick] (p) -- node[above] {map} ++(2.0,0);
\draw[-{Latex}, thick, dashed] ($(p)+(0,-0.45)$) -- node[below] {pullback / inverse image} ++(2.0,0);
\end{tikzpicture}
\end{center}

This matters because diagrams should not become decoration. They are a second
notation. The point is to help you see maps, domains, targets, identifications,
and failure points before the symbols tighten the idea.

\begin{marginnotebox}
When a diagram appears, ask what each region, arrow, and color means formally.
If you cannot translate the picture back into definitions, the picture is only
a sketch.
\end{marginnotebox}

\subsection{Recurring Motifs}

The book also uses small recurring marks as cognitive landmarks.

\begin{longtable}{p{0.14\linewidth}p{0.28\linewidth}p{0.48\linewidth}}
\toprule
Mark & Name & Meaning \\
\midrule
\motifcompression & compression arrow & examples or repeated moves are being compressed into a definition, theorem, or structure \\
\motiflocalglobal & local/global marker & a local fact is being assembled into a global conclusion, or a global condition controls local behavior \\
\motifhiddenquantifier & hidden quantifier & a phrase is hiding ``for every,'' ``there exists,'' or a dependency of choices \\
\motifstrategy & strategy fork & the proof could go several ways; pause and choose deliberately \\
\bottomrule
\end{longtable}

These marks are not new mathematics. They are memory handles. When the material
gets abstract, familiar landmarks keep you oriented.

\subsection{Abstraction as Compression}

One of the central themes of the whole reader is this:

\begin{quote}
Abstraction is compression with memory.
\end{quote}

At first you learn examples. Then a definition compresses many examples into a
single object. A theorem compresses many repeated arguments into one reusable
tool. A structure, such as group or topology, compresses many different worlds
that obey the same rules. A structure-preserving map compresses the idea of
``same kind of behavior'' across different objects.

This is why mathematics becomes more abstract as it matures. The goal is not to
make things vague. The goal is to store more meaning in fewer, more powerful
concepts.

\begin{aestheticbox}
Good abstraction feels beautiful because it suddenly makes distant examples
look like relatives. Modular arithmetic, symmetries of a square, and invertible
matrices are different on the surface, but group theory lets one proof speak to
all of them.
\end{aestheticbox}

\begin{recoverybox}
Try this with no symbols. Name three things in your life that are different on
the surface but follow the same pattern. A playlist, a task queue, and a proof
outline all have order. A map, an API, and a function all send inputs somewhere.
This is the emotional beginning of abstraction: noticing repeated shape.
\end{recoverybox}

\subsection{The Emotional Weather of Learning Proofs}

There are several normal states in learning rigorous mathematics:

\begin{itemize}
\item \textbf{Notation fatigue:} the symbols are individually clear, but the page feels noisy.
\item \textbf{Delayed understanding:} a definition makes sense only after three examples and one failed proof.
\item \textbf{Abstraction vertigo:} the text has moved from familiar objects to a structure that contains them.
\item \textbf{Proof amnesia:} you understood a proof yesterday but cannot reconstruct it today.
\item \textbf{Semantic blur:} you can read every word but cannot yet say what the theorem is really doing.
\end{itemize}

None of these are signs that you are not mathematical. They are signs that your
mind is building new compression machinery.

\begin{focusbox}
The cure for abstraction vertigo is not speed. It is returning to examples,
non-examples, and the exact definition until the abstract word has weight again.
\end{focusbox}

\subsection{How To Slow Down Productively}

There is an unhelpful way to slow down: stare at a page until the symbols become familiar-looking. There is also a productive way: interrogate the sentence.

When you meet a claim, ask:

\begin{itemize}
\item What objects are being discussed?
\item What universe do those objects live in?
\item Is the sentence true or false once the objects are fixed?
\item Which variables are already known, and which are to be chosen?
\item Is the claim universal, existential, conditional, or some mixture?
\item What would count as a proof?
\item What would count as a counterexample?
\end{itemize}

These questions are not extra decorations. They are how proofs start.

For instance:
\[
\forall x\in\R,\exists y\in\R,\quad y>x.
\]
This statement is not asking for a single magic number $y$. It gives you $x$ first, then allows you to answer with a $y$ depending on $x$. A proof can choose $y=x+1$. If you miss the order of the quantifiers, the statement feels mysterious. If you see the order, it is almost mechanical.

\begin{walkbox}
When you walk away from a definition, ask: what are the inputs, what are the outputs, and what promises are being made?
\end{walkbox}

\subsection{The Change From Examples To Reasons}

Examples are good. Examples are necessary. Examples are not proofs.

Suppose you notice:
\[
1=1^2,\qquad 1+3=2^2,\qquad 1+3+5=3^2,\qquad 1+3+5+7=4^2.
\]
It is natural to guess:
\[
1+3+5+\cdots+(2n-1)=n^2.
\]
Checking the first hundred cases gives evidence. Checking the first million cases gives more evidence. A proof gives something else: a reason the pattern cannot fail.

One proof is visual. Start with a $1\times1$ square. To build a $2\times2$ square, add an L-shaped border of $3$ dots. To build a $3\times3$ square, add an L-shaped border of $5$ dots. To build an $n\times n$ square, the last border has $2n-1$ dots. So the first $n$ odd numbers build an $n\times n$ square.

Another proof uses induction. Another uses the arithmetic series formula. Another counts the number of pairs $(i,j)$ with $1\leq i,j\leq n$ by grouping them according to the larger of $i$ and $j$. The theorem is the same, but the proofs reveal different structure.

\begin{intuitionbox}
A computation samples the landscape. A proof explains why the landscape has that shape.
\end{intuitionbox}

This is a central shift. In proof-based mathematics, an answer without a reason is unfinished. A reason without an answer may still be valuable, because it may explain why a whole family of answers behaves as it does.

\subsection{What Counts As Understanding}

For the early foundations, understanding is not a private feeling of familiarity. It has observable signs.

You understand a statement when you can:

\begin{itemize}
\item classify it as a statement, predicate, expression, command, or question;
\item identify its variables and their roles;
\item state its domain;
\item negate it accurately;
\item say what a proof would need to do;
\item say what a disproof would need to exhibit;
\item produce small examples and edge cases.
\end{itemize}

That is a lot. It is supposed to be. Mathematical reading is active.

Take the statement:
\[
\forall n\in\Z,\quad n^2+n\text{ is even}.
\]
Understanding it includes knowing that $n$ is arbitrary, the domain is $\Z$, the claim is universal, and a counterexample would be one integer $n$ for which $n^2+n$ is odd. It also includes seeing at least two proof routes:
\[
n^2+n=n(n+1),
\]
the product of consecutive integers, one of which is even; or cases, where $n=2k$ or $n=2k+1$.

You do not need both proofs every time. But seeing multiple routes makes the terrain more real.

\subsection{Formal Does Not Mean Inhuman}

Some people talk as if rigor is a punishment: first you enjoy mathematics, then rigor arrives and takes the toys away. That is not the right picture.

Rigor is what lets us keep promises. It lets us distinguish a theorem from a pattern, a definition from a vibe, a proof from a convincing diagram. The goal is not to make mathematics sterile. The goal is to make it honest enough that other people can trust it, build on it, and find exactly where it fails.

\begin{historybox}
Greek mathematics, especially Euclid's \emph{Elements}, made proof a public object: not just a private conviction, but a chain of reasons others could inspect. Centuries later, calculus forced mathematicians to sharpen the language again. Infinitesimals, limits, continuity, and convergence all worked brilliantly in practice before their foundations were made precise. The history of rigor is not the history of people being fussy. It is the history of powerful ideas becoming powerful enough to need better grammar.
\end{historybox}

This history matters because you will sometimes feel that a definition is too careful. ``Why do we have to say domain and codomain?'' ``Why does the order of quantifiers matter?'' ``Why distinguish a statement from a predicate?'' The answer is usually: because later, if we do not, something breaks.

\subsection{The Four Views}

For each concept in this part, we will keep returning to four views.

\textbf{The informal view.} What is the plain-language idea?

\textbf{The formal view.} What is the exact statement?

\textbf{The proof view.} What kind of action does the statement demand?

\textbf{The debugging view.} How do beginners commonly misuse it?

For example, the universal quantifier $\forall$ has these views:

\begin{itemize}
\item Informal: every object in the domain has the property.
\item Formal: $\forall x\in D,\ P(x)$.
\item Proof: let $x\in D$ be arbitrary and prove $P(x)$.
\item Debugging: do not prove it by checking a few favorite values of $x$.
\end{itemize}

The existential quantifier $\exists$ has parallel views:

\begin{itemize}
\item Informal: at least one object in the domain has the property.
\item Formal: $\exists x\in D$ such that $P(x)$.
\item Proof: produce or justify a witness.
\item Debugging: do not act as if the witness works uniformly unless the statement requires that.
\end{itemize}

The same idea viewed four ways becomes much harder to confuse.

\subsection{A Software Analogy: Types, Scope, and Contracts}

If you have programmed, logic will feel less alien when you compare it to types and scope.

A variable with no stated domain is like a parameter with no type. The expression might still be meaningful in someone's head, but the machine cannot safely use it.

The predicate
\[
P(x): x^2<10
\]
depends on the type of $x$. If $x\in\N$, the satisfying values are limited. If $x\in\Z$, there are negative values too. If $x\in\R$, there is an interval. If $x$ is a matrix, the expression $x^2<10$ may not even mean anything until we define an order.

A bound variable is like a local variable. In
\[
\sum_{i=1}^{n} i^2,
\]
the $i$ is local to the summation. If the next line says ``let $i$ be an integer,'' that is a new use of the letter unless the context says otherwise. The letter is not sacred; the role is.

A theorem is like a contract. It says: if you give me objects satisfying these hypotheses, I promise this conclusion. A proof is the implementation of the contract.

\begin{trapbox}
Do not confuse familiar notation with well-defined notation. A symbol you have seen before can become meaningless if the domain changes.
\end{trapbox}

\subsection{A Physical Analogy: Instruments and Specifications}

Imagine a lab instrument with a specification:

\begin{quote}
For every target temperature between $10^\circ$C and $40^\circ$C, there exists a dial setting that reaches that temperature within one minute.
\end{quote}

That has quantifier structure:
\[
\forall T\in[10,40],\exists s\in S,\quad R(s,T).
\]
The setting may depend on the target. If the target is $12^\circ$C, use one setting. If the target is $35^\circ$C, use another.

Now compare:
\[
\exists s\in S,\forall T\in[10,40],\quad R(s,T).
\]
This says one fixed setting reaches every target temperature. That is a much stronger and probably false claim.

Physical systems make the order of choices obvious. Mathematics often hides the same issue behind compact notation.

\subsection{Beginner-To-Rigor Reading Routine}

When you face a theorem, use this routine.

\begin{enumerate}
\item Strip the theorem down to its logical skeleton.
\item Name the domain of every variable.
\item Mark which variables are arbitrary and which are chosen.
\item Translate the conclusion into the action needed to prove it.
\item Try one example.
\item Try one non-example or boundary case.
\item Write the first line of the proof before thinking about the clever part.
\end{enumerate}

The first line is often forced. For
\[
\forall x\in A,\exists y\in B,\quad P(x,y),
\]
the proof probably begins:

\begin{quote}
Let $x\in A$ be arbitrary. We must find $y\in B$ such that $P(x,y)$.
\end{quote}

That line does not solve the problem, but it puts the problem in the correct shape.

\begin{formalbox}
A proof is not a transcript of how you discovered the idea. It is a structured explanation showing why the conclusion follows from the assumptions.
\end{formalbox}

\subsection{Proof Lab: One Claim, Three Readings}

\begin{labbox}
Study the claim:
\[
\forall n\in\Z,\quad n^2-n\text{ is even}.
\]

\textbf{Reading 1: examples.}
\[
3^2-3=6,\qquad (-4)^2-(-4)=20,\qquad 0^2-0=0.
\]
The examples support the claim, but they do not prove it.

\textbf{Reading 2: structure.}
\[
n^2-n=n(n-1).
\]
The numbers $n$ and $n-1$ are consecutive integers, so one is even. Therefore their product is even.

\textbf{Reading 3: proof skeleton.}
Let $n\in\Z$ be arbitrary. Since $n$ and $n-1$ are consecutive integers, one of them is even. Hence $n(n-1)$ is even. Therefore $n^2-n$ is even. Since $n$ was arbitrary, the claim holds for every integer $n$.
\end{labbox}

Notice the last sentence. It is not decoration. It closes the universal proof by reminding the reader that no special property of $n$ was used.

\subsection{Non-Examples Of Proof}

The following are not proofs of
\[
\forall n\in\Z,\quad n^2-n\text{ is even}.
\]

\textbf{Several examples.}
\[
1^2-1=0,\quad 2^2-2=2,\quad 3^2-3=6.
\]
Good evidence, not proof.

\textbf{An unsupported pattern claim.}

\begin{quote}
It keeps alternating, so it must be even.
\end{quote}

What alternates? Why must it continue? What theorem justifies the pattern?

\textbf{A circular explanation.}

\begin{quote}
$n^2-n$ is even because it is always divisible by $2$.
\end{quote}

That only restates the definition of evenness. It does not show divisibility by $2$.

\textbf{A false general principle.}

\begin{quote}
Since $n^2$ and $n$ have the same parity, their difference is even.
\end{quote}

This one is actually a valid route if you prove the parity statement first. Without proof, it is a hidden lemma.

\begin{trapbox}
When a proof says ``clearly,'' ask what theorem, definition, or earlier result makes it clear. Sometimes ``clearly'' is fine. Sometimes it is where the missing proof is hiding.
\end{trapbox}

\subsection{Exercises}

\begin{enumerate}
\item For the claim $\forall n\in\Z,\ n^2+n$ is even, write one example, one proof idea, and one full proof.
\item Give a computational check for $1+3+\cdots+(2n-1)=n^2$ when $n=6$. Then explain why that check is not a proof.
\item Find a formula from calculus or physics that you know. State one assumption under which it is valid and one situation where it may fail.
\item Write the first line of a proof of: $\forall x\in\R,\exists y\in\R$ such that $x+y=0$.
\item Write the first line of a disproof of: $\exists y\in\R,\forall x\in\R$ such that $x+y=0$.
\item Explain in ordinary language why the order of choices matters in a thermostat or control-system specification.
\end{enumerate}

\begin{walkbox}
Think of a theorem as a promise. Who is allowed to choose the inputs? Who must produce the outputs? What would count as breaking the promise?
\end{walkbox}

\section{Mathematical Sentences}

Before proof, there is grammar. Not grammar in the sense of commas and adjectives, but mathematical grammar: what kind of object is this line of writing?

Some lines assert something. Some merely name an object. Some ask a question. Some issue an instruction. Some contain variables and become true or false only after the variables receive values. If you do not know which kind of line you are reading, proof becomes blurry before it even starts.

\begin{formalbox}
A \emph{statement} is a sentence that is either true or false.
\end{formalbox}

Examples:

\begin{itemize}
\item ``$7$ is prime'' is a true statement.
\item ``$9$ is prime'' is a false statement.
\item ``There are infinitely many prime numbers'' is a true statement.
\item ``Every continuous function is differentiable'' is a false statement.
\item ``There are infinitely many twin primes'' is a statement whose truth is not currently known.
\end{itemize}

The last example is important. A sentence can be a statement even if humans do not know whether it is true. Truth value and known truth value are different ideas.

\subsection{Statements Are Claims, Not Decorations}

The expression
\[
x^2+1
\]
is not a statement. It does not claim anything. It is an expression, a mathematical object waiting to be evaluated or used.

The line
\[
x^2+1=10
\]
looks more sentence-like, but by itself it is not yet a statement unless $x$ already has a fixed meaning. If $x=3$, it is true. If $x=0$, it is false. If $x$ is unspecified, the line is a predicate: a condition that may become true or false after $x$ is supplied.

The command

\begin{quote}
Solve $x^2+1=10$.
\end{quote}

is not a statement. It is an instruction.

The question

\begin{quote}
Is $x^2+1=10$?
\end{quote}

is not a statement. It asks for one.

\begin{focusbox}
The first proof-reading filter is simple: does this sentence have a truth value yet?
\end{focusbox}

\subsection{Predicates}

\begin{formalbox}
A \emph{predicate} is a sentence with one or more variables that becomes a statement when the variables are assigned values from a specified domain.
\end{formalbox}

For example,
\[
P(n):\quad n\text{ is prime}
\]
is a predicate. If $n=7$, $P(n)$ is true. If $n=9$, $P(n)$ is false. If $n=1$, $P(n)$ is false under the standard definition of prime.

Another predicate:
\[
Q(x):\quad x^2<2.
\]
This predicate changes personality depending on the domain.

Over $\N=\{1,2,3,\ldots\}$, it is true only for $x=1$.

Over $\Z$, it is true for $x=-1,0,1$ if $0$ is included in the integers, as usual.

Over $\R$, it is true exactly when
\[
-\sqrt2<x<\sqrt2.
\]

Same symbols. Different domains. Different mathematical situation.

\begin{trapbox}
Never let a variable float without a domain. A variable without a domain is a character without a setting.
\end{trapbox}

\subsection{Expressions, Equations, Inequalities, And Sentences}

It helps to separate four categories.

\textbf{Expression.} A mathematical object or recipe for one:
\[
x^2+1,\qquad \sin t,\qquad A\vec v,\qquad \sum_{k=1}^{n}k.
\]
Expressions can be simplified, evaluated, transformed, or compared. They are not true or false by themselves.

\textbf{Equation or inequality with free variables.} A condition:
\[
x^2+1=10,\qquad x<y,\qquad \|v\|=1.
\]
These may become statements after variables are assigned or quantified.

\textbf{Statement.} A complete claim:
\[
\forall x\in\R,\quad x^2+1>0.
\]
This has a truth value.

\textbf{Instruction.} A command to do something:
\[
\text{Find all }x\in\R\text{ such that }x^2+1=10.
\]
The instruction itself is not true or false, although a proposed answer can be correct or incorrect.

In textbooks, these categories are often mixed in quick succession. A problem might say:

\begin{quote}
Let $f(x)=x^2+1$. Solve $f(x)=10$. Prove that $f(x)>0$ for all $x\in\R$.
\end{quote}

The first sentence defines notation. The second gives an instruction. The third states a proof task. You need to hear the shift.

\subsection{Open Sentences And Closed Sentences}

An \emph{open sentence} has a free variable. A \emph{closed sentence} has no free variables. Closed mathematical sentences are candidates for being statements.

Open:
\[
x^2=4.
\]

Closed by assignment:
\[
3^2=4.
\]
This is false.

Closed by quantification:
\[
\exists x\in\R,\quad x^2=4.
\]
This is true.

Closed another way:
\[
\forall x\in\R,\quad x^2=4.
\]
This is false.

The same open sentence can be closed in different ways, producing different statements.

\begin{intuitionbox}
An open sentence is like a form with blanks. A quantifier tells you how the blanks are to be filled: by every possible value, by at least one value, or by values chosen in a particular order.
\end{intuitionbox}

\subsection{Truth Is Relative To Interpretation}

The sentence
\[
2+2=4
\]
feels absolutely true, and in ordinary arithmetic it is. But even here, mathematics is more careful than habit.

In modular arithmetic modulo $3$,
\[
2+2\equiv 1\pmod 3.
\]
That does not make ordinary arithmetic false. It means the symbols are being interpreted in a different structure.

Similarly,
\[
x^2\geq0
\]
is true for all real numbers $x$. But if $x$ is complex, the symbol $\geq$ is not naturally defined in the same order-compatible way. The sentence is not automatically false; it may be ill-formed unless an order has been specified.

In linear algebra, the expression $AB=BA$ may be true for some pairs of matrices and false for others. In a commutative ring it is true for all elements. In a noncommutative ring it may fail. The same visual pattern of symbols does not settle the meaning by itself.

\begin{historybox}
The nineteenth and twentieth centuries forced mathematicians to become more explicit about interpretation. Non-Euclidean geometry showed that familiar words like ``line'' and ``parallel'' could satisfy different axiom systems. Abstract algebra showed that operations resembling addition or multiplication could live in many structures. Logic responded by separating syntax, the written form, from semantics, the interpretation that gives it meaning.
\end{historybox}

\subsection{A Small Truth-Value Table For Sentences}

Classify the following lines.

\[
5\in\Z
\]
This is a true statement.

\[
\pi\in\Q
\]
This is a false statement.

\[
x\in\Q
\]
This is a predicate if $x$ is not already fixed.

\[
\text{Let }x\in\Q.
\]
This is not a statement. It introduces a variable under a condition.

\[
\text{Show that }\sqrt2\notin\Q.
\]
This is an instruction.

\[
\sqrt2\notin\Q.
\]
This is a true statement.

These distinctions are easy when displayed slowly. In real proofs, they are compressed. A line like ``Let $x\in A$'' may be followed immediately by ``Then $x\in B$,'' and you need to know that the first line sets context while the second asserts a claim under that context.

\subsection{Mathematical Sentences In Proofs}

Proofs are written in ordinary language mixed with symbols. That does not mean every sentence in a proof is a theorem. Proof text contains:

\begin{itemize}
\item assumptions,
\item variable introductions,
\item definitions being unfolded,
\item intermediate claims,
\item justifications,
\item conclusions.
\end{itemize}

Consider this short proof:

\begin{quote}
Let $n\in\Z$ be even. Then $n=2k$ for some $k\in\Z$. Hence $n^2=4k^2=2(2k^2)$, and since $2k^2\in\Z$, $n^2$ is even.
\end{quote}

The sentence ``Let $n\in\Z$ be even'' introduces a fixed but arbitrary integer satisfying a hypothesis. The sentence ``$n=2k$ for some $k\in\Z$'' is a claim justified by the definition of even. The equation $n^2=4k^2=2(2k^2)$ is a chain of equalities. The final phrase closes the definition.

Reading a proof well means hearing which role each line plays.

\subsection{Claims And Their Social Names}

Mathematics uses several names for proved statements.

\begin{itemize}
\item A \emph{theorem} is usually a major result.
\item A \emph{proposition} is a result of moderate importance.
\item A \emph{lemma} is a helper result used to prove something else.
\item A \emph{corollary} follows quickly from an earlier result.
\item A \emph{claim} is often a local statement inside a proof.
\end{itemize}

These names are cultural labels, not rigid biological species. One author's lemma may be another author's theorem. The logical status is the same: each is a statement intended to be proved from accepted assumptions, definitions, and previous results.

\subsection{Non-Statements That Look Mathematical}

Here are common non-statements.

\textbf{A bare expression:}
\[
\frac{x^2-1}{x-1}.
\]
It may simplify to $x+1$ when $x\neq1$, but by itself it is not true or false.

\textbf{A definition command:}
\[
\text{Define }f:\R\to\R\text{ by }f(x)=x^2.
\]
This sets notation.

\textbf{A problem instruction:}
\[
\text{Determine whether }A\subseteq B.
\]
This asks you to produce a statement, either $A\subseteq B$ or $A\not\subseteq B$, with justification.

\textbf{A dangling implication:}
\[
\text{If }x>0.
\]
This is incomplete. An implication needs a conclusion.

\textbf{A type declaration:}
\[
x\in\R.
\]
This can be a statement if $x$ has already been named. It can also be a variable declaration if it appears as ``Let $x\in\R$.'' Context matters.

\begin{trapbox}
The line $x\in\R$ is not automatically one thing in all contexts. In a proof, ``Let $x\in\R$'' introduces a variable. In a list of assumptions, $x\in\R$ is a condition. In a completed argument about a previously defined object $x$, it may be a claim.
\end{trapbox}

\subsection{Worked Example: Sorting The Sentence Types}

Classify each line.

\[
4+6=10.
\]
Statement, true.

\[
n+6=10.
\]
Predicate, unless $n$ has already been assigned.

\[
n+6.
\]
Expression.

\[
\text{Find }n\in\Z\text{ such that }n+6=10.
\]
Instruction.

\[
\forall n\in\Z,\quad n+6=10.
\]
Statement, false.

\[
\exists n\in\Z,\quad n+6=10.
\]
Statement, true.

\[
\text{Is there an }n\in\Z\text{ such that }n+6=10?
\]
Question.

The open condition $n+6=10$ can be converted into different statements by different quantifiers. This is why classifying the sentence comes before proving it.

\subsection{Worked Example: Domain Changes Meaning}

Consider:
\[
P(x):\quad x^2=2.
\]

Over $\N$, no $x$ satisfies $P(x)$.

Over $\Z$, no $x$ satisfies $P(x)$.

Over $\Q$, no $x$ satisfies $P(x)$, though proving that requires the irrationality of $\sqrt2$.

Over $\R$, two values satisfy it: $x=\sqrt2$ and $x=-\sqrt2$.

Over $\C$, the same two real values satisfy it, but now they are viewed as complex numbers with imaginary part $0$.

So the statement
\[
\exists x,\quad x^2=2
\]
is incomplete until the domain of $x$ is known. Over $\Q$ it is false. Over $\R$ it is true.

\begin{focusbox}
Whenever you see ``there exists $x$,'' train yourself to ask: exists where?
\end{focusbox}

\subsection{Software Analogy: Boolean Functions And Type Errors}

A predicate resembles a Boolean-valued function:

\[
\texttt{isPrime}(n)
\]
returns true or false when $n$ is a suitable input. But if you call
\[
\texttt{isPrime}(\texttt{"hello"})
\]
you have not discovered a false prime. You have made a type error.

Mathematics has similar issues. The predicate
\[
x<5
\]
is natural for real numbers. It is not automatically meaningful for vectors, matrices, functions, or sets. You might define an order on some of these objects, but that definition is extra structure. Without it, the sentence is not ready.

This is one reason pure mathematics is careful about domains. The domain is not bureaucratic paperwork. It is part of the meaning.

\subsection{Physical Analogy: Measuring Devices}

A measuring device produces meaningful readings only under a protocol. A thermometer reading ``37'' means different things depending on Celsius, Fahrenheit, calibration, and where the sensor is placed. The number alone is not the full statement.

Similarly,
\[
x=37
\]
means little until $x$ is known. Is $x$ a temperature, an index, a voltage, a dimension, a solution to an equation, an element of a finite field? The symbol can be reused, but meaning comes from context.

Mathematical notation is compact because context carries weight. Beginners often blame themselves for not understanding a symbol, when the real problem is that they have not identified the surrounding sentence.

\subsection{Proof Lab: From Predicate To Theorem}

\begin{labbox}
Start with the predicate:
\[
P(n):\quad n^2+n\text{ is even}.
\]

\textbf{Version 1: universal theorem.}
\[
\forall n\in\Z,\quad P(n).
\]
This says every integer satisfies the predicate. A proof must handle an arbitrary integer.

\textbf{Version 2: existential theorem.}
\[
\exists n\in\Z\quad P(n).
\]
This says at least one integer satisfies the predicate. A proof may simply choose $n=1$, since $1^2+1=2$ is even.

\textbf{Version 3: false universal theorem over a changed predicate.}
\[
\forall n\in\Z,\quad n^2+n+1\text{ is even}.
\]
This is false. A counterexample is $n=0$, giving $1$.

\textbf{Lesson.}
The predicate is the raw condition. The quantifier turns it into a theorem-shaped claim.
\end{labbox}

\subsection{Common Traps}

\begin{trapbox}
\textbf{Trap 1: Treating an equation with variables as automatically true.}
The line $x^2=9$ is not true for all $x$. It is a condition satisfied by some values and not others.

\textbf{Trap 2: Forgetting the domain.}
The statement ``$x^2=2$ has a solution'' is true over $\R$ and false over $\Q$.

\textbf{Trap 3: Confusing a command with a claim.}
``Find $x$'' is not something you prove. A proposed answer is what gets checked.

\textbf{Trap 4: Thinking unknown means neither true nor false.}
The twin prime conjecture is a statement even though its truth is not known.
\end{trapbox}

\subsection{Exercises}

Classify each as statement, predicate, expression, command, question, or context-setting phrase. If it is a predicate, give two possible domains and explain how the truth set changes.

\begin{enumerate}
\item $3+4=7$.
\item $x+4=7$.
\item $x+4$.
\item Find all $x\in\R$ such that $x+4=7$.
\item Is every integer a rational number?
\item Every integer is a rational number.
\item $A\subseteq B$.
\item Let $A$ and $B$ be sets.
\item $\sin^2 t+\cos^2 t=1$.
\item $\sin^2 t+\cos^2 t$.
\item $\forall t\in\R,\ \sin^2 t+\cos^2 t=1$.
\item There exists an integer $n$ such that $n^2=12$.
\item There exists a real number $x$ such that $x^2=12$.
\item Define $f:\R\to\R$ by $f(x)=x^2$.
\item $f$ is injective.
\end{enumerate}

For items involving $A,B,$ or $f$, state what extra context is needed before the line has a definite truth value.

\begin{walkbox}
During ordinary reading, pause at a displayed formula and ask: is this an object, a condition, a claim, or an instruction?
\end{walkbox}

\section{Variables Have Roles}

Beginners often treat variables as generic letters. In proofs, variables have jobs. Confusing those jobs is one of the fastest ways to get lost.

The same letter can play different roles in different sentences. The role is not determined by the shape of the letter. It is determined by the logic around it.

\begin{focusbox}
When a variable appears, ask: who introduced it, what domain does it live in, and am I allowed to choose it?
\end{focusbox}

\subsection{Fixed But Unknown}

Suppose a proof begins:

\begin{quote}
Let $n\in\Z$.
\end{quote}

Now $n$ is fixed for the duration of the local argument. We may not know whether $n$ is $5$, $-12$, or $0$, but it is one particular integer. It should not quietly change values halfway through the proof.

This is similar to a function parameter in software. If a function is called with an input, that input may be unknown to the programmer while writing the function, but during one execution it has one value.

\[
\texttt{def f(n):}
\]

Inside the function body, $n$ is not all integers at once. It is the input value for that call. A proof often treats a variable similarly: fixed but unspecified.

For example:

\begin{quote}
Let $n\in\Z$. Suppose $n$ is even.
\end{quote}

Now $n$ is a fixed integer with an additional property. From evenness we may write:
\[
n=2k\quad\text{for some }k\in\Z.
\]
The $k$ is not arbitrary. It is chosen as a witness to the evenness of this fixed $n$.

\subsection{Arbitrary}

When a proof says:

\begin{quote}
Let $n\in\Z$ be arbitrary.
\end{quote}

it means we are not allowed to use any special feature of $n$ beyond membership in $\Z$ and whatever hypotheses are stated. If the argument works for such an arbitrary $n$, it works for all integers.

This is the standard opening move for proving a universal statement:
\[
\forall n\in\Z,\quad P(n).
\]

Proof skeleton:

\begin{quote}
Let $n\in\Z$ be arbitrary. We will show $P(n)$.
\end{quote}

The word ``arbitrary'' is not magic, but it is a warning label. It says: do not secretly choose $n=2$ because that is convenient. Do not assume $n$ is positive unless positivity is part of the hypothesis. Do not assume $n$ is nonzero unless nonzero is part of the hypothesis.

\begin{trapbox}
To prove a statement about every integer, you may not prove it only for your favorite integer. ``Let $n=4$'' proves something about $4$, not about all $n\in\Z$.
\end{trapbox}

\subsection{Chosen As A Witness}

Existential statements ask for witnesses.

\[
\exists n\in\Z\quad n^2=49.
\]

A witness is an object that makes the existential claim true. Here $n=7$ is a witness, and so is $n=-7$.

To prove:
\[
\exists x\in\R\quad x^2-5x+6=0,
\]
you can choose $x=2$ and verify:
\[
2^2-5\cdot2+6=4-10+6=0.
\]

That is enough. You do not need to find every solution unless the problem asks for all solutions. Existential proof requires at least one witness, not a census.

Sometimes the witness is built from earlier variables. To prove:
\[
\forall x\in\R,\exists y\in\R\quad x+y=0,
\]
you cannot choose one $y$ before seeing $x$. After $x$ is given, choose:
\[
y=-x.
\]
Then $x+y=x+(-x)=0$.

\begin{intuitionbox}
A witness is evidence with a name. For an existential claim, one correctly verified witness is enough.
\end{intuitionbox}

\subsection{Bound Variables}

A variable is \emph{bound} when its meaning is controlled by a quantifier or a local operator.

In
\[
\forall x\in\R,\quad x^2\geq0,
\]
the variable $x$ is bound by $\forall x\in\R$.

In
\[
\exists k\in\Z,\quad n=2k,
\]
the variable $k$ is bound by $\exists k\in\Z$.

In
\[
\sum_{i=1}^{n} i^2,
\]
the variable $i$ is bound by the summation. It is local to the sum.

In
\[
\int_a^b f(t)\,dt,
\]
the variable $t$ is bound by the integral. It is a dummy variable. The following expressions are the same:
\[
\int_a^b f(t)\,dt
\qquad\text{and}\qquad
\int_a^b f(u)\,du.
\]

The letter changed; the role did not.

\begin{formalbox}
Bound variables can usually be renamed, as long as the renaming does not collide with another variable and change the meaning.
\end{formalbox}

\subsection{Free Variables}

A variable is \emph{free} when it is not bound by a quantifier or local operator and has not been fixed by surrounding context.

In
\[
x^2\geq0,
\]
$x$ is free unless earlier text has introduced it.

In
\[
\forall x\in\R,\quad x<y,
\]
$x$ is bound, but $y$ is free unless context defines it. This is not a complete statement by itself. It becomes meaningful if, for example, $y$ has already been fixed as a real number. Then it says every real number is less than that fixed $y$, which is false.

Free variables are not bad. They are how we write predicates, definitions, and local claims. But they require context.

\subsection{Dummy Variables}

Some bound variables are called dummy variables because their names carry no lasting meaning.

\[
\sum_{i=1}^{n} i
=
\sum_{j=1}^{n} j.
\]

The $i$ and $j$ are dummy variables. They are placeholders used by the summation notation.

Likewise:
\[
\{x\in\R:x^2<1\}
=
\{t\in\R:t^2<1\}.
\]

Both describe the open interval $(-1,1)$.

But not every renaming is safe. Compare:
\[
\forall x\in\R,\exists y\in\R,\quad y>x.
\]
If we foolishly rename $y$ to $x$ inside the same scope, we get:
\[
\forall x\in\R,\exists x\in\R,\quad x>x,
\]
which is broken and misleading. The inner $x$ shadows the outer $x$, and the condition $x>x$ is never true. The renaming caused a collision.

\begin{trapbox}
Renaming bound variables is safe only when it preserves scope. Avoid reusing a letter if the old and new variables need to be compared.
\end{trapbox}

\subsection{Parameters}

A parameter is a fixed object that controls a family of statements or objects.

For example:
\[
f_a(x)=ax
\]
defines a family of functions indexed by $a$. Here $x$ is the input variable, while $a$ is a parameter. If $a=2$, the function is $f_2(x)=2x$. If $a=-1$, the function is $f_{-1}(x)=-x$.

In the statement:
\[
\forall a\in\R,\forall x\in\R,\quad f_a(x)=ax,
\]
both $a$ and $x$ are bound by universal quantifiers. But conceptually, $a$ chooses a member of the family and $x$ chooses an input to that member.

Parameters are common in analysis:
\[
\lim_{x\to a} f(x)=L.
\]
Here $a$ and $L$ are often fixed while $x$ varies toward $a$. If you blur these roles, epsilon-delta proofs become painful.

\subsection{Constants And Named Objects}

A constant is a fixed named object. Some constants are universal, like $\pi$ or $e$. Others are local:

\begin{quote}
Let $c=\sqrt2$.
\end{quote}

Now $c$ has a fixed meaning in the current context.

But a letter is not permanently a constant just because it was once used that way. In one proof $c$ may be a fixed real number. In another proof $c$ may be a witness chosen after an epsilon is given. The local text decides.

Mathematical letters are reusable tools. Do not attach too much personality to them.

\subsection{Variables In Definitions}

Definitions often use variables in a special way. Consider:

\[
n\text{ is even}\quad\Longleftrightarrow\quad \exists k\in\Z\text{ such that }n=2k.
\]

Here $n$ is the object being defined. The variable $k$ is a witness variable inside the definition. To prove $n$ is even, you must produce such a $k$. To use the fact that $n$ is even, you are allowed to introduce such a $k$.

Example:

\begin{quote}
Suppose $n$ is even. Then there exists $k\in\Z$ such that $n=2k$.
\end{quote}

At this point, you often write:

\begin{quote}
Choose $k\in\Z$ such that $n=2k$.
\end{quote}

This $k$ depends on $n$. If $n=10$, one possible $k$ is $5$. If $n=-6$, one possible $k$ is $-3$.

\begin{focusbox}
Definitions are machines with variables inside. To use a definition well, know which variables are inputs and which variables must be supplied as witnesses.
\end{focusbox}

\subsection{Scope}

Scope is the region of text where a variable has a particular meaning.

In a proof:

\begin{quote}
Let $x\in A$. Since $A\subseteq B$, we have $x\in B$.
\end{quote}

The variable $x$ is in scope for those sentences. If the proof later starts a new paragraph:

\begin{quote}
Now let $x\in C$.
\end{quote}

the writer may be reusing the same letter for a new object. Good writers avoid confusing reuse, but textbooks sometimes rely on context.

In nested quantifiers, scope is even more important:
\[
\forall x\in A,\exists y\in B,\forall z\in C,\quad R(x,y,z).
\]
The $y$ may depend on $x$, because $y$ is chosen after $x$. The $z$ is chosen after both $x$ and $y$, but because it is universal, the claim must work for all $z$.

\subsection{Dependency}

Dependency is one of the most important ideas hidden in variable roles.

Compare:
\[
\forall x\in\R,\exists y\in\R,\quad y=x^2
\]
and
\[
\exists y\in\R,\forall x\in\R,\quad y=x^2.
\]

The first is true. Given $x$, choose $y=x^2$. The witness $y$ depends on $x$.

The second is false. It asks for one fixed real number $y$ equal to every square $x^2$. No single number can be both $0^2$ and $2^2$.

Dependency is not always written explicitly. Sometimes authors write:
\[
\forall x\in X,\exists y_x\in Y,\quad P(x,y_x).
\]
The subscript $y_x$ is a reminder that the chosen $y$ may depend on $x$. But the quantifier order already said that. The notation only makes the dependency visible.

\begin{intuitionbox}
Quantifier order is a dependency diagram in disguise.
\end{intuitionbox}

\subsection{Worked Parsing: A Simple Nested Claim}

Consider:
\[
\forall x\in\R,\exists y\in\R\text{ such that }y>x.
\]

Read the quantifiers in order from left to right:

\begin{enumerate}
\item The challenger gives an arbitrary real number $x$.
\item After seeing $x$, we must produce a real number $y$.
\item The produced $y$ must satisfy $y>x$.
\end{enumerate}

A proof:

\begin{quote}
Let $x\in\R$ be arbitrary. Choose $y=x+1$. Then $y\in\R$, and since $x+1>x$, we have $y>x$. Therefore for every real $x$ there exists a real $y$ such that $y>x$.
\end{quote}

Variable roles:

\begin{itemize}
\item $x$ is arbitrary.
\item $y$ is a witness.
\item $y$ depends on $x$.
\end{itemize}

Now reverse the quantifiers:
\[
\exists y\in\R,\forall x\in\R\text{ such that }y>x.
\]

Now $y$ must be chosen first. It cannot depend on $x$. The claim asks for a real number strictly greater than every real number. But taking $x=y$ would require $y>y$, which is impossible; hence no such $y$ exists and the statement is false.

\subsection{Worked Parsing: Epsilon Language}

Analysis uses variable roles intensely. The statement
\[
\lim_{n\to\infty}a_n=L
\]
means:
\[
\forall \eps>0,\exists N\in\N,\forall n\in\N,\quad n\geq N\Rightarrow |a_n-L|<\eps.
\]

Roles:

\begin{itemize}
\item $\eps$ is chosen by the challenger. It is an arbitrary positive tolerance.
\item $N$ is chosen in response to $\eps$.
\item $n$ is arbitrary after $N$ is chosen, subject to $n\geq N$.
\item $L$ is usually a fixed proposed limit.
\end{itemize}

The dependency is:
\[
N=N(\eps).
\]
The $N$ may depend on $\eps$. It may not depend on $n$, because $N$ is chosen before the final universal $n$.

\begin{trapbox}
In an epsilon proof, choosing $N$ in terms of $n$ is a scope error. The proof must find one threshold $N$ that works for all later $n\geq N$.
\end{trapbox}

\subsection{Worked Parsing: Set Inclusion}

The statement
\[
A\subseteq B
\]
means:
\[
\forall x,\quad x\in A\Rightarrow x\in B.
\]

In a proof of $A\subseteq B$, the variable $x$ is arbitrary:

\begin{quote}
Let $x\in A$ be arbitrary. We must show $x\in B$.
\end{quote}

This is why set-inclusion proofs often feel repetitive. They are rehearsing the logic of the subset symbol.

To disprove $A\subseteq B$, you need a witness:
\[
\exists x\quad x\in A\text{ and }x\notin B.
\]

The same variable $x$ is arbitrary in the proof of inclusion and a witness in the disproof. The role changes with the task.

\subsection{Physical Analogy: Knobs And Readings}

In a lab setup, some quantities are controlled, some are measured, and some are fixed by the apparatus.

Suppose a statement says:

\begin{quote}
For every input voltage $V$, there exists a resistor setting $R$ such that the current is below $1$ amp.
\end{quote}

Here $V$ is arbitrary; $R$ is a response. The setting $R$ may depend on $V$.

If the statement says:

\begin{quote}
There exists a resistor setting $R$ such that for every input voltage $V$, the current is below $1$ amp.
\end{quote}

then $R$ is fixed before $V$ varies. That is a different engineering promise.

Mathematical variables behave the same way. Some are knobs. Some are readings. Some are universal disturbances. Some are chosen controls.

\subsection{Software Analogy: Scope Bugs}

Consider pseudocode:

\[
\texttt{for x in X:}
\]
\[
\qquad\texttt{choose y depending on x}
\]

This corresponds to:
\[
\forall x\in X,\exists y\in Y,\quad P(x,y).
\]

Now compare:

\[
\texttt{choose y}
\]
\[
\texttt{for x in X:}
\]
\[
\qquad\texttt{check P(x,y)}
\]

This corresponds to:
\[
\exists y\in Y,\forall x\in X,\quad P(x,y).
\]

The code is different. The mathematics is different. The order of lines controls what data is available.

\subsection{Common Variable Traps}

\begin{trapbox}
\textbf{Trap 1: Changing an arbitrary variable into a special one.}
If $n$ is arbitrary, do not later say ``take $n=2$'' unless you are switching to an example.

\textbf{Trap 2: Choosing a witness too early.}
In $\forall x\exists y$, $y$ can depend on $x$. In $\exists y\forall x$, it cannot.

\textbf{Trap 3: Reusing a bound variable outside its scope.}
The $i$ in $\sum_{i=1}^{n}i^2$ has no life after the summation ends.

\textbf{Trap 4: Treating a parameter as a variable of motion.}
In $\lim_{x\to a}f(x)=L$, the variable $x$ moves; $a$ and $L$ are fixed in the statement.
\end{trapbox}

\subsection{Proof Lab: Repairing Variable Misuse}

\begin{labbox}
\textbf{Bad proof attempt.}
Claim: $\forall n\in\Z,\ n^2+n$ is even.

``Let $n=3$. Then $n^2+n=12$, which is even. Therefore $n^2+n$ is even for all integers.''

\textbf{Diagnosis.}
The proof chooses a special $n$. It proves the statement for $n=3$, not for arbitrary $n$.

\textbf{Repaired proof.}
Let $n\in\Z$ be arbitrary. Then
\[
n^2+n=n(n+1).
\]
The integers $n$ and $n+1$ are consecutive, so one of them is even. Therefore their product is even. Hence $n^2+n$ is even. Since $n$ was arbitrary, the result holds for all $n\in\Z$.
\end{labbox}

\subsection{Proof Lab: Witnesses Depending On Earlier Choices}

\begin{labbox}
Claim:
\[
\forall a\in\R,\exists b\in\R,\quad a+b=5.
\]

\textbf{Proof.}
Let $a\in\R$ be arbitrary. Choose $b=5-a$. Since $a\in\R$, we have $5-a\in\R$. Also
\[
a+b=a+(5-a)=5.
\]
Thus for every real number $a$, there exists a real number $b$ such that $a+b=5$.

\textbf{Role check.}
The variable $a$ is arbitrary. The variable $b$ is a witness. The witness depends on $a$.
\end{labbox}

\subsection{Exercises}

For each statement, identify the role of every variable: arbitrary, fixed parameter, witness, bound dummy variable, or free variable needing context.

\begin{enumerate}
\item $\forall n\in\Z,\exists k\in\Z$ such that $n^2+n=2k$.
\item $\exists k\in\Z,\forall n\in\Z$ such that $n^2+n=2k$.
\item $\sum_{i=1}^{n}i^2$.
\item $\int_0^1 x^n\,dx$.
\item $\forall \eps>0,\exists \delta>0,\forall x\in\R,\ |x-a|<\delta\Rightarrow |f(x)-f(a)|<\eps$.
\item $A\subseteq B$.
\item $\forall x\in A,\exists y\in B,\forall z\in C,\ R(x,y,z)$.
\item $\exists y\in B,\forall x\in A,\forall z\in C,\ R(x,y,z)$.
\end{enumerate}

Then answer:

\begin{enumerate}
\item In which statements may the existential variable depend on an earlier universal variable?
\item In which statements must one object work for all later choices?
\item Which variables can be safely renamed?
\item Which variables would cause confusion if renamed carelessly?
\end{enumerate}

\begin{walkbox}
When reading a proof, track variables like characters entering and leaving a scene. Who is on stage? Who has already been chosen? Who is allowed to respond to whom?
\end{walkbox}

\section{Quantifiers}

The symbols $\forall$ and $\exists$ are small, but they control the shape of entire proofs.

\[
\forall=\text{for every},\qquad \exists=\text{there exists}.
\]

They are not merely shorthand. They are instructions for a mathematical game of choices.

\begin{focusbox}
To understand a quantified statement, read it as an order of play: who chooses first, who responds, and what must be true after the choices are made?
\end{focusbox}

\subsection{Universal Statements}

A universal statement has the form:
\[
\forall x\in D,\quad P(x).
\]

It says every object in the domain $D$ satisfies property $P$.

To prove it, use the universal proof pattern:

\begin{quote}
Let $x\in D$ be arbitrary. We must show $P(x)$.
\end{quote}

To disprove it, find one counterexample:

\begin{quote}
Find $x\in D$ such that $P(x)$ is false.
\end{quote}

Example:
\[
\forall n\in\Z,\quad n^2\geq0.
\]
True, because the square of any integer is nonnegative.

False universal statement:
\[
\forall n\in\Z,\quad n>0.
\]
A counterexample is $n=0$, or $n=-1$.

The counterexample does not need to be dramatic. One valid failure is enough.

\subsection{Existential Statements}

An existential statement has the form:
\[
\exists x\in D\quad P(x).
\]

It says at least one object in $D$ satisfies $P$.

To prove it, produce a witness:

\begin{quote}
Choose $x=\text{something in }D$. Verify $P(x)$.
\end{quote}

To disprove it, prove a universal negative:

\begin{quote}
Let $x\in D$ be arbitrary. Show $P(x)$ is false.
\end{quote}

Example:
\[
\exists n\in\Z,\quad n^2=16.
\]
True. A witness is $n=4$.

False existential statement:
\[
\exists n\in\Z,\quad n^2=2.
\]
False, because no integer has square $2$.

Notice the asymmetry:

\begin{itemize}
\item A universal statement needs all cases for proof, one case for disproof.
\item An existential statement needs one case for proof, all cases for disproof.
\end{itemize}

\subsection{The Quantifier Game}

Think of $\forall$ as a challenger and $\exists$ as a responder.

For
\[
\forall x\in\R,\exists y\in\R,\quad x+y=0,
\]
the game goes:

\begin{enumerate}
\item Challenger chooses $x\in\R$.
\item You choose $y\in\R$.
\item You must make $x+y=0$ true.
\end{enumerate}

Winning strategy:
\[
y=-x.
\]

For
\[
\exists y\in\R,\forall x\in\R,\quad x+y=0,
\]
the game goes:

\begin{enumerate}
\item You choose one $y\in\R$.
\item Challenger chooses $x\in\R$.
\item The equation $x+y=0$ must be true.
\end{enumerate}

There is no winning strategy. After you choose $y$, the challenger can choose $x=1-y$, giving
\[
x+y=1.
\]

The two statements use the same equation but have different quantifier order. One is true; the other is false.

\begin{intuitionbox}
Quantifiers are not seasoning sprinkled on a formula. They are part of the claim.
\end{intuitionbox}

\subsection{Order Matters}

Compare:
\[
\forall x\in\R,\exists y\in\R,\quad y>x
\]
and
\[
\exists y\in\R,\forall x\in\R,\quad y>x.
\]

First statement: true. Given $x$, choose $y=x+1$.

Second statement: false. It asks for a real number larger than every real number.

Compare:
\[
\forall x\in\R,\exists y\in\R,\quad y^2=x
\]
and
\[
\exists y\in\R,\forall x\in\R,\quad y^2=x.
\]

The first is false over $\R$, because negative $x$ values have no real square roots. The second is also false, because one fixed square cannot equal every real number.

Now change the domain:
\[
\forall x\in[0,\infty),\exists y\in\R,\quad y^2=x.
\]
This is true. Given $x\geq0$, choose $y=\sqrt{x}$.

\begin{focusbox}
Quantifier order and domain are two of the main controls on mathematical meaning.
\end{focusbox}

\subsection{Restricted Quantifiers}

The notation
\[
\forall x\in D,\quad P(x)
\]
is shorthand for:
\[
\forall x,\quad x\in D\Rightarrow P(x).
\]

The notation
\[
\exists x\in D\quad P(x)
\]
is shorthand for:
\[
\exists x,\quad x\in D\text{ and }P(x).
\]

This difference matters when negating statements later. Universal restrictions pair naturally with implication; existential restrictions pair naturally with ``and.''

Example:
\[
\forall x\in\R,\quad x^2\geq0
\]
means:
\[
\forall x,\quad x\in\R\Rightarrow x^2\geq0.
\]

Example:
\[
\exists x\in\R\quad x^2=2
\]
means:
\[
\exists x,\quad x\in\R\text{ and }x^2=2.
\]

The restricted notation is easier to read, but the expanded form reveals the logic.

\subsection{Vacuous Truth}

Universal statements over empty domains are true.

\[
\forall x\in\emptyset,\quad P(x)
\]
is true, no matter what $P$ is.

This feels strange at first. The reason is counterexamples. To make a universal statement false, you need an object in the domain that fails the property. If the domain is empty, there is no such object.

Example:
\[
\forall x\in\emptyset,\quad x^2=17
\]
is true.

But:
\[
\exists x\in\emptyset,\quad x^2=17
\]
is false.

A physical analogy: ``Every device in the empty box is broken'' and ``Every device in the empty box works'' are both vacuously true if there are no devices in the box. The existential claim ``There is a working device in the empty box'' is false.

\begin{trapbox}
Vacuous truth is not a loophole. It is what keeps universal statements logically consistent when domains are empty.
\end{trapbox}

\subsection{Uniqueness}

Sometimes existence is not enough. We want exactly one object.

The notation
\[
\exists!x\in D\quad P(x)
\]
means:

\begin{quote}
There exists exactly one $x\in D$ such that $P(x)$.
\end{quote}

This contains two claims:

\begin{enumerate}
\item Existence: at least one such $x$ exists.
\item Uniqueness: any two such objects must be equal.
\end{enumerate}

Formal expansion:
\[
\exists x\in D\left(P(x)\text{ and }\forall y\in D,\ P(y)\Rightarrow y=x\right).
\]

Example:
\[
\exists!x\in\R\quad x+3=10.
\]
True, because $x=7$ works, and if $y+3=10$, then $y=7$.

Example:
\[
\exists!x\in\R\quad x^2=4.
\]
False, because both $x=2$ and $x=-2$ work.

\begin{labbox}
To prove uniqueness, do not merely display one solution. After existence, suppose $u$ and $v$ are both solutions, then prove $u=v$.
\end{labbox}

\subsection{At Least, At Most, Exactly}

Quantifiers also express counting claims.

\textbf{At least one:}
\[
\exists x\in D\quad P(x).
\]

\textbf{At most one:}
\[
\forall x,y\in D,\quad (P(x)\text{ and }P(y))\Rightarrow x=y.
\]

\textbf{Exactly one:}
\[
\exists x\in D\quad P(x)
\]
and
\[
\forall x,y\in D,\quad (P(x)\text{ and }P(y))\Rightarrow x=y.
\]

Example:

The equation $x^2=0$ has exactly one real solution. Existence: $x=0$ works. At most one: if $u^2=0$ and $v^2=0$, then $u=0$ and $v=0$, so $u=v$.

The equation $x^2=1$ has at least one real solution, but not at most one, because $1$ and $-1$ both work.

\subsection{Quantifiers And Implication}

Many mathematical statements have the form:
\[
\forall x\in D,\quad P(x)\Rightarrow Q(x).
\]

This says every object in $D$ that satisfies $P$ also satisfies $Q$.

Example:
\[
\forall n\in\Z,\quad n\text{ is divisible by }4\Rightarrow n\text{ is even}.
\]

To prove it:

\begin{quote}
Let $n\in\Z$ be arbitrary. Suppose $n$ is divisible by $4$. Show $n$ is even.
\end{quote}

You do not need to prove every integer is divisible by $4$. You assume the hypothesis for an arbitrary integer and prove the conclusion under that assumption.

To disprove it, you need:
\[
\exists n\in\Z\quad n\text{ is divisible by }4\text{ and }n\text{ is not even}.
\]
No such $n$ exists, so the original statement is true.

\subsection{Common Speech Versus Mathematical Speech}

Ordinary language is flexible. Mathematical language is less forgiving.

``Everyone has someone they admire'' usually means:
\[
\forall p\in P,\exists q\in P,\quad A(p,q).
\]
Each person may admire a different person.

``There is someone everyone admires'' means:
\[
\exists q\in P,\forall p\in P,\quad A(p,q).
\]
One person is admired by all.

These are very different social claims.

Likewise:

\begin{quote}
Every problem has a method that solves it.
\end{quote}

means:
\[
\forall p,\exists m,\quad m\text{ solves }p.
\]

But:

\begin{quote}
There is a method that solves every problem.
\end{quote}

means:
\[
\exists m,\forall p,\quad m\text{ solves }p.
\]

The second is much stronger.

\subsection{History And Connections}

\begin{historybox}
Quantifiers became explicit relatively late in the history of mathematics. Mathematicians had always used words like ``all'' and ``some,'' but modern symbolic logic, especially in the work of Frege, Peano, Russell, Hilbert, and others, made quantifier structure an object of study. This mattered because mathematics was becoming more abstract. Once you study arbitrary sets, functions, spaces, and structures, you need a language that says exactly which objects are being varied and which are fixed.
\end{historybox}

Quantifiers later appear everywhere.

In analysis:
\[
\forall \eps>0,\exists \delta>0,\forall x,\quad |x-a|<\delta\Rightarrow |f(x)-f(a)|<\eps.
\]

In algebra:
\[
\forall a\in G,\exists b\in G,\quad ab=ba=e.
\]

In topology:
\[
\forall x\in U,\exists r>0,\quad B(x,r)\subseteq U.
\]

In computer science:
\[
\forall \text{ inputs }i,\exists \text{ output }o,\quad \text{Spec}(i,o).
\]

Quantifiers are not a chapter you leave behind. They are the operating system.

\subsection{Worked Example: Proving A Universal-Existential Statement}

Claim:
\[
\forall n\in\Z,\exists m\in\Z,\quad n+m=0.
\]

Proof:

\begin{quote}
Let $n\in\Z$ be arbitrary. Choose $m=-n$. Since $n\in\Z$, we have $-n\in\Z$, so $m\in\Z$. Also
\[
n+m=n+(-n)=0.
\]
Therefore for every integer $n$, there exists an integer $m$ such that $n+m=0$.
\end{quote}

What to notice:

\begin{itemize}
\item The proof starts with an arbitrary $n$ because the first quantifier is universal.
\item The witness $m$ is chosen after $n$.
\item The proof verifies both membership $m\in\Z$ and the required equation.
\end{itemize}

\subsection{Worked Example: Disproving A Reversed Statement}

Claim:
\[
\exists m\in\Z,\forall n\in\Z,\quad n+m=0.
\]

This is false.

Disproof:

\begin{quote}
Suppose, for contradiction, that there exists $m\in\Z$ such that for all $n\in\Z$, $n+m=0$. Since the statement holds for all integers $n$, it holds for $n=0$ and $n=1$. Taking $n=0$ gives $m=0$. Taking $n=1$ gives $1+m=0$, so $m=-1$. Thus $m=0$ and $m=-1$, a contradiction. Therefore no such integer $m$ exists.
\end{quote}

You can also give a direct disproof:

\begin{quote}
Let $m\in\Z$ be arbitrary. Choose $n=1-m$. Then $n\in\Z$, but $n+m=1\neq0$. So this $m$ does not work. Since $m$ was arbitrary, no integer $m$ works for all $n$.
\end{quote}

The direct disproof is especially instructive: to defeat a proposed universal property of $m$, choose a counterexample $n$ depending on $m$.

\subsection{Worked Example: A Three-Quantifier Statement}

Consider:
\[
\forall a\in\R,\forall b\in\R,\exists x\in\R,\quad x>a\text{ and }x>b.
\]

Plain language:

\begin{quote}
Given any two real numbers, there is a real number larger than both.
\end{quote}

Proof:

\begin{quote}
Let $a,b\in\R$ be arbitrary. Choose
\[
x=\max\{a,b\}+1.
\]
Then $x\in\R$, and $x>a$ and $x>b$. Therefore for any real numbers $a$ and $b$, there exists a real number larger than both.
\end{quote}

If you do not want to use $\max$, use:
\[
x=a^2+b^2+|a|+|b|+1.
\]
That is overkill, but it works. Proofs do not always need the prettiest witness; they need a correct witness.

\subsection{Worked Example: A False Three-Quantifier Statement}

Consider:
\[
\exists x\in\R,\forall a\in\R,\forall b\in\R,\quad x>a\text{ and }x>b.
\]

This says there is one real number larger than all pairs of real numbers. In particular, it would be larger than every real number. False.

Disproof:

\begin{quote}
Let $x\in\R$ be arbitrary. Choose $a=x+1$ and $b=0$. Then $a\in\R$ and $b\in\R$, but $x>a$ is false. Therefore this $x$ does not satisfy the required property. Since $x$ was arbitrary, no such real $x$ exists.
\end{quote}

Again, the counterexample may depend on the proposed witness. That is legal because in disproving $\exists x\forall a\forall b$, we start by considering an arbitrary candidate $x$ and then show it fails.

\subsection{Quantifiers In Definitions}

Definitions often hide quantifiers.

\textbf{Even integer.}
\[
n\text{ is even}\quad\Longleftrightarrow\quad \exists k\in\Z,\ n=2k.
\]

\textbf{Subset.}
\[
A\subseteq B\quad\Longleftrightarrow\quad \forall x,\ x\in A\Rightarrow x\in B.
\]

\textbf{Function equality.}
If $f,g:A\to B$, then
\[
f=g\quad\Longleftrightarrow\quad \forall x\in A,\ f(x)=g(x).
\]

\textbf{Injective function.}
\[
f:A\to B\text{ is injective}
\quad\Longleftrightarrow\quad
\forall x_1,x_2\in A,\ f(x_1)=f(x_2)\Rightarrow x_1=x_2.
\]

\textbf{Surjective function.}
\[
f:A\to B\text{ is surjective}
\quad\Longleftrightarrow\quad
\forall y\in B,\exists x\in A,\ f(x)=y.
\]

Surjectivity is a quantifier statement. The output target $y$ is arbitrary; the input $x$ is a witness that may depend on $y$.

\begin{trapbox}
Many proof mistakes are definition mistakes: the writer remembers the slogan but not the quantifiers inside the definition.
\end{trapbox}

\subsection{Negating Quantifiers Briefly}

Although negation gets its own deeper treatment later, quantifiers already demand the basic rules:

\[
\neg(\forall x\in D,\ P(x))
\quad\Longleftrightarrow\quad
\exists x\in D\text{ such that }\neg P(x).
\]

\[
\neg(\exists x\in D,\ P(x))
\quad\Longleftrightarrow\quad
\forall x\in D,\ \neg P(x).
\]

Not every integer is positive means there exists an integer that is not positive.

There is no integer whose square is $2$ means every integer has square not equal to $2$.

For nested quantifiers, flip each quantifier as the negation passes through:

\[
\neg\left(\forall x\in A,\exists y\in B,\ P(x,y)\right)
\]
becomes
\[
\exists x\in A,\forall y\in B,\ \neg P(x,y).
\]

Plain language:

\begin{quote}
It is not true that every $x$ has some good $y$.
\end{quote}

means:

\begin{quote}
There is a particular $x$ for which every $y$ fails.
\end{quote}

\subsection{Proof Lab: Translating A Definition}

\begin{labbox}
Definition:
\[
f:A\to B\text{ is surjective}
\quad\Longleftrightarrow\quad
\forall b\in B,\exists a\in A,\ f(a)=b.
\]

\textbf{Proof skeleton for surjectivity.}
Let $b\in B$ be arbitrary. We must find $a\in A$ such that $f(a)=b$.

\textbf{Disproof skeleton.}
Find $b\in B$ such that for every $a\in A$, $f(a)\neq b$.

\textbf{Example.}
Let $f:\R\to\R$ be $f(x)=x^3$. To prove $f$ is surjective, let $b\in\R$ be arbitrary and choose $a=\sqrt[3]{b}$. Then $f(a)=a^3=b$.

\textbf{Non-example.}
Let $g:\R\to\R$ be $g(x)=x^2$. It is not surjective, because $b=-1\in\R$, and for every $a\in\R$, $a^2\neq-1$.
\end{labbox}

\subsection{Proof Lab: Quantifier Order In Limits}

\begin{labbox}
The sequence statement
\[
\lim_{n\to\infty}a_n=L
\]
means:
\[
\forall \eps>0,\exists N\in\N,\forall n\in\N,\quad n\geq N\Rightarrow |a_n-L|<\eps.
\]

\textbf{Correct dependency.}
$N$ may depend on $\eps$.

\textbf{Incorrect dependency.}
$N$ may not depend on $n$, because the statement requires the same $N$ to work for all $n\geq N$.

\textbf{Example.}
For $a_n=1/n$ and $L=0$, given $\eps>0$, choose $N$ so large that $N>1/\eps$. Then for every $n\geq N$,
\[
\left|\frac1n-0\right|=\frac1n\leq \frac1N<\eps.
\]
\end{labbox}

\subsection{Common Quantifier Traps}

\begin{trapbox}
\textbf{Trap 1: Reversing quantifiers casually.}
$\forall x\exists y$ and $\exists y\forall x$ usually say different things.

\textbf{Trap 2: Proving an existential when the theorem is universal.}
One example cannot prove an all-claim.

\textbf{Trap 3: Proving a universal when only existence was needed.}
This is not logically wrong, but it may waste effort. If one witness is enough, one witness is enough.

\textbf{Trap 4: Forgetting membership.}
When you choose a witness, verify it lies in the required domain.

\textbf{Trap 5: Letting later variables influence earlier choices.}
In $\forall\eps\exists N\forall n$, the $N$ cannot depend on $n$.
\end{trapbox}

\subsection{A Quantifier Checklist}

Before proving a quantified statement, write:

\begin{enumerate}
\item The quantifiers in order.
\item The role of each variable.
\item What you are allowed to choose.
\item What must be arbitrary.
\item What the final condition is.
\item What a counterexample would look like.
\end{enumerate}

For:
\[
\forall x\in X,\exists y\in Y,\forall z\in Z,\quad P(x,y,z),
\]
the checklist says:

\begin{itemize}
\item $x$ is arbitrary.
\item $y$ is chosen after $x$ and may depend on $x$.
\item $z$ is arbitrary after $y$ is chosen.
\item The proof must show $P(x,y,z)$ for every $z\in Z$.
\item A disproof needs one $x\in X$ such that every $y\in Y$ fails for at least one $z\in Z$.
\end{itemize}

The negation makes the last item precise:
\[
\exists x\in X,\forall y\in Y,\exists z\in Z,\quad \neg P(x,y,z).
\]

\subsection{Visualizing Quantifier Flow}

Quantifiers are easier to read when you stop seeing them as a string of symbols and start seeing them as a flow of responsibility. The order tells you who moves first, who responds, and what information is available at each step.

\begin{center}
\begin{tikzpicture}[
  node distance=1.35cm,
  every node/.style={font=\small},
  qnode/.style={draw, rounded corners, align=center, minimum width=3.2cm, minimum height=0.85cm, fill=blue!5},
  enode/.style={draw, rounded corners, align=center, minimum width=3.2cm, minimum height=0.85cm, fill=green!6},
  arrow/.style={-{Latex}, thick}
]
\node[qnode] (x) {Challenger chooses\\$x\in X$};
\node[enode, right=of x] (y) {You choose\\$y=y(x)\in Y$};
\node[qnode, right=of y] (z) {Challenger chooses\\$z\in Z$};
\node[enode, right=of z] (p) {You prove\\$P(x,y,z)$};
\draw[arrow] (x) -- (y);
\draw[arrow] (y) -- (z);
\draw[arrow] (z) -- (p);
\node[below=0.55cm of y, align=center, text width=7.5cm] {For $\forall x\in X,\exists y\in Y,\forall z\in Z,\ P(x,y,z)$, the witness $y$ may depend on $x$, but it must be chosen before $z$ is revealed.};
\end{tikzpicture}
\end{center}

That one sentence, ``chosen before $z$ is revealed,'' is often the difference between a valid proof and a proof that accidentally cheats.

\begin{trapbox}
False intuition: ``If I can choose a good object eventually, the existential quantifier is satisfied.'' Not always. You must choose it at the point where the quantifier appears. In $\exists y\,\forall x$, the $y$ is chosen before seeing $x$. In $\forall x\,\exists y$, the $y$ may respond to $x$.
\end{trapbox}

\subsection{Symbolic Execution for Quantifiers}

Software engineers are used to tracing a program with inputs, local variables, and assertions. Quantified statements can be traced the same way.

For the statement
\[
\forall x\in X,\exists y\in Y,\forall z\in Z,\quad P(x,y,z),
\]
the execution trace is:

\begin{enumerate}
\item Input arrives: an arbitrary $x\in X$.
\item Your code must compute or specify a $y\in Y$. This $y$ may use $x$.
\item Another input arrives: an arbitrary $z\in Z$.
\item Your assertion must pass: $P(x,y,z)$.
\end{enumerate}

The negation is the failing test case:
\[
\exists x\in X,\forall y\in Y,\exists z\in Z,\quad \neg P(x,y,z).
\]
It says there is a particular $x$ that breaks every possible response $y$ by producing some bad $z$.

\begin{labbox}
Proof anatomy drill. For each quantified statement, write the execution trace and the failing test case.
\begin{enumerate}
\item $\forall n\in\Z,\exists m\in\Z,\ n< m$.
\item $\forall \eps>0,\exists N\in\N,\forall n\geq N,\ |a_n-L|<\eps$.
\item $\exists \delta>0,\forall x,y\in\R,\ |x-y|<\delta\Rightarrow |f(x)-f(y)|<1$.
\end{enumerate}
\end{labbox}

\subsection{Exercises}

Decide whether each statement is true or false. Prove it or give a counterexample.

\begin{enumerate}
\item $\forall x\in\R,\exists y\in\R,\quad x<y$.
\item $\exists y\in\R,\forall x\in\R,\quad x<y$.
\item $\forall n\in\Z,\exists m\in\Z,\quad n<m$.
\item $\exists m\in\Z,\forall n\in\Z,\quad n<m$.
\item $\forall x\in\R,\exists y\in\R,\quad x+y=10$.
\item $\exists y\in\R,\forall x\in\R,\quad x+y=10$.
\item $\forall x\in[0,\infty),\exists y\in\R,\quad y^2=x$.
\item $\forall x\in\R,\exists y\in\R,\quad y^2=x$.
\item $\exists x\in\R,\forall y\in\R,\quad x+y=y$.
\item $\forall y\in\R,\exists x\in\R,\quad x+y=y$.
\end{enumerate}

Translate each into symbols.

\begin{enumerate}
\item Every real number has an additive inverse.
\item Some real number is larger than every integer.
\item For every integer, there is a larger integer.
\item There is a smallest natural number.
\item Every nonempty finite set has an element.
\item There is a function from $\R$ to $\R$ that sends every real number to $0$.
\end{enumerate}

Negate each statement.

\begin{enumerate}
\item $\forall x\in A,\ P(x)$.
\item $\exists x\in A,\ P(x)$.
\item $\forall x\in A,\exists y\in B,\ P(x,y)$.
\item $\exists y\in B,\forall x\in A,\ P(x,y)$.
\item $\forall \eps>0,\exists N\in\N,\forall n\geq N,\ |a_n-L|<\eps$.
\end{enumerate}

\begin{walkbox}
Pick a sentence from ordinary life containing ``everyone,'' ``someone,'' or ``always.'' Translate it into quantifiers. Then reverse the quantifiers and notice how the meaning changes.
\end{walkbox}
